\scanpage{54}
b) Sur $S_D\deleted{\unsure{=S'}}$, on a un rev. [étale], qui doit prolonger le rev. étale qu'on a déjà sur $S_{2D}\deleted{\unsure{=S'_2}}$

On est donc amené \deleted{(oubliant les pts de $V(D)=S-S_D$)} à traiter ce pb de prolongmt, en se rappellant les données $(L',D')$ sur lesquelles il faut ``\unsure{modeler|recoller}'' le prolongmt, \deleted{\unsure{lorsqu'on}} Un premier invariant important est le torseur $F_{\deleted{\unsure{q}}}$\inserted{\deleted{\unsure{$=F|S'$}}} (noté simplement $F$) sous $\check{\mathbb V}(L)$, trivialisé sur $S'_2$. En fait, on a mieux, savoir une section \inserted{$\underline b$} de $2F$, donc une section $\inserted{b_0=}\frac12\underline b$ de $F|S'_2$ \deleted{\unsure{$\subset V|S$}} dont le produit par $f=2$ se prolonge, et est muni d'un prolongmt. En termes de l'inclusion
\[i:S'_f\longleftarrow S'\qquad(f=2\cdot1_{S'})\]
on peut dire que l'on a
\[\deleted{\unsure{\Gamma(S',}}\ \underline L\subset\tfrac1f\underline L\subset i_*(L|S_f)\]
en supposant $f=2$ section régulière de $\unsure{\mathcal O_S|\mathcal O_{S'}}$, et $F$ est un torseur sous $\underline L$ tel que le torseur correspondant sous $\frac1f\underline L$ soit trivialisé. \inserted{La donnée d'} \unsure{Un|un} tel torseur \deleted{et de} équivaut à la donnée d'une section de $\frac1f\underline L/\underline L\simeq\underline L/f\underline L$, \deleted{\unsure{savoir}} obtenue ainsi : on choisit \inserted{(loc !)} une trivialisation \deleted{locale} $s_0$ de $F$, on écrit $b_0$ comme section de $L|S_f$ de façon précise :
\[b_0=s_{\unsure{*|0}}+\tfrac1f x\qquad(x\in\Gamma(S,L))\]
% archive sheet number: 53
